% !TeX root = ../master.tex

\chapter{Experiments and Results}
\label{ch:experiments}

%%%%%%%%%%%%%%%%%%%%%%%%%%%%%%%%%%%%%%%%%%%%%%%
%% Part 1: Zero-Shot Robustness (RQ1)        %%
%%%%%%%%%%%%%%%%%%%%%%%%%%%%%%%%%%%%%%%%%%%%%%%

\section{Zero-Shot Robustness Across Model Scales}
\label{sec:zero_shot}

% Part 1: RQ1 -- Do zero-shot LLMs rely on argument structure?

\subsection{Experimental Setup}
\label{sec:zs_setup}

% - All 5 models, all 10 datasets, balanced samples
% - Context: dataset-specific argument definition (C1 baseline)
% - 3 manipulations per sample (original, Feger, shuffle)
% - Parallel classification via ThreadPoolExecutor

\subsection{Results}
\label{sec:zs_results}

% - Table: Macro F1 per model x dataset for M0, M1, M2
% - Table: Delta_feger and Delta_shuffle per model x dataset
% - Overall means across datasets per model
% - Key findings from existing experiments:
%   - GPT-4.1 baseline results (full_30, baseline_30 experiments)
%   - Mistral-Small-24B results (baseline_30, full_30)
%   - Llama-8B and Llama-70B results
%   - Mistral-7B results
% - Preliminary: pilot on ABSTRCT showed delta_feger=-0.10, delta_shuffle=-0.29

\subsection{Model Scale Analysis}
\label{sec:zs_scale}

% - How does delta vary with model size? (7B -> 8B -> 24B -> 70B -> frontier)
% - Is there a threshold where models become robust/sensitive?
% - Do larger models show larger delta (more reliance on structure)?
% - Or smaller delta (more robust to perturbation)?
% - Compare across model families: Mistral vs. Llama at similar sizes

\subsection{Per-Dataset Analysis}
\label{sec:zs_per_dataset}

% - Which datasets are hardest/easiest for zero-shot LLMs?
% - Dataset-specific delta patterns: do some datasets have more exploitable shortcuts?
% - Relate to dataset characteristics: domain, annotation criteria, class balance
% - Compare to encoder results from Feger et al. where available

\subsection{Comparison to Encoder Baselines}
\label{sec:zs_comparison}

% - Direct comparison: LLM deltas vs. encoder delta <= 0.02
% - 5x larger drop for Feger manipulation (pilot: delta=-0.10 vs. 0.02)
% - Shuffle as new diagnostic: encoders not tested, but order-invariance suggests small delta
% - Interpretation: LLMs use argument structure that encoders ignore
% - Caveats: zero-shot vs. fine-tuned comparison, different absolute F1 levels

\subsection{Contamination Analysis}
\label{sec:zs_contamination}

% - Results from contamination_test_outputs/contamination_test_20260205.json
% - Per-model, per-dataset ROUGE-L and exact match rates
% - Flag any model-dataset pairs with ROUGE-L > 0.5
% - Temporal analysis: performance vs. dataset publication year
% - Impact on interpretation of zero-shot results

%%%%%%%%%%%%%%%%%%%%%%%%%%%%%%%%%%%%%%%%%%%%%%%
%% Part 2: Context Utilization (RQ2)         %%
%%%%%%%%%%%%%%%%%%%%%%%%%%%%%%%%%%%%%%%%%%%%%%%

\section{Context Utilization}
\label{sec:context}

% Part 2: RQ2 -- Does context improve zero-shot argument identification?

\subsection{Experimental Setup}
\label{sec:ctx_setup}

% - Context conditions C0--C4:
%   - C0: generic argument definition ("An argument is a statement that...")
%   - C1: dataset-specific argument definition (extracted from paper)
%   - C2: C1 + annotation guidelines (4 datasets only)
%   - C3: C1 + document context (2 preceding sentences, 6 datasets)
%   - C4: C1 + guidelines + document context (4 datasets with full availability)
% - Existing experiments map to conditions:
%   - no_context_30 = C0 (simple definition)
%   - baseline_30 = C1 (definition only)
%   - guidelines_30 = C2 (definition + guidelines)
%   - full_30 = C4 (definition + guidelines + document context)

\subsection{Results: Definition vs. No Definition}
\label{sec:ctx_definition}

% - C0 vs. C1: generic vs. dataset-specific definition
% - Pilot results: average F1 jumps from 0.58 to 0.63 across all 10 datasets
% - Non-uniform improvement: AFS and ACQUA benefit most, some datasets negative
% - Interpretation: value depends on divergence between generic and actual annotation criteria

\subsection{Results: Annotation Guidelines}
\label{sec:ctx_guidelines}

% - C1 vs. C2: does adding annotation guidelines help?
% - Only 4 datasets have guidelines: ABSTRCT, ARGUMINSCI, PE, USELEC
% - Compare guidelines_30 vs. baseline_30 experiments
% - Cross-guideline experiment results (cross_guideline_outputs):
%   - What happens when you give dataset A the guidelines of dataset B?
%   - Tests whether guidelines provide dataset-specific vs. general argument knowledge

\subsection{Results: Document Context}
\label{sec:ctx_document}

% - C1 vs. C3: does adding preceding sentences help?
% - 6 datasets with document context: ABSTRCT, ARGUMINSCI, FINARG, PE, SCIARK, USELEC
% - Hypothesis: surrounding text disambiguates borderline cases
% - Compare full_30 (has doc context) vs. guidelines_30 (no doc context)

\subsection{Results: Combined Context}
\label{sec:ctx_combined}

% - C4: full context (definition + guidelines + document context)
% - Does combining all sources yield additive gains?
% - Or diminishing returns / interference between context sources?
% - Full_30 experiment results across models

\subsection{Context and Manipulation Sensitivity}
\label{sec:ctx_manipulation}

% - Does context change the model's reliance on sentence-internal structure?
% - Compare delta_feger and delta_shuffle across context conditions
% - Hypothesis: with more external context, model may rely less on sentence structure
% - Or: context helps the model focus on the right structural cues

\subsection{Per-Model Analysis}
\label{sec:ctx_per_model}

% - Do smaller models benefit more from context (compensating for weaker reasoning)?
% - Or do larger models leverage context better?
% - Model x context interaction effects
% - Practical implications: when is context worth the added prompt length?

%%%%%%%%%%%%%%%%%%%%%%%%%%%%%%%%%%%%%%%%%%%%%%%
%% Part 3: Fine-Tuning Effects (RQ3)         %%
%%%%%%%%%%%%%%%%%%%%%%%%%%%%%%%%%%%%%%%%%%%%%%%

\section{Fine-Tuning Effects}
\label{sec:fine_tuning}

% Part 3: RQ3 -- Does fine-tuning reintroduce shortcut learning?

\subsection{Experimental Setup}
\label{sec:ft_setup}

% - Model: best-performing open-weight model from Part 1 + LoRA adapters
% - Training data: full GAIC training set (all 10 datasets)
% - Input: same prompt structure as Parts 1--2, full context condition (C4)
% - LoRA configuration: rank, alpha, target modules, learning rate
% - Training details: epochs, batch size, optimizer, early stopping criteria

\subsection{Fine-Tuning Results}
\label{sec:ft_results}

% - F1 on standard test sets (original sentences)
% - Comparison: zero-shot vs. fine-tuned absolute performance
% - Per-dataset improvements or degradations

\subsection{Manipulation Robustness After Fine-Tuning}
\label{sec:ft_robustness}

% - F1 on Feger-manipulated test sets
% - F1 on shuffled test sets
% - Delta comparison: fine-tuned vs. zero-shot
% - Key interpretation table:
%   - delta <= 0.05: encoder-like, shortcut learning is training-dependent
%   - delta 0.05--0.20: partial resistance, decoders partially resist
%   - delta >= 0.20: robustness preserved, decoders fundamentally differ

\subsection{Cross-Dataset Transfer}
\label{sec:ft_transfer}

% - Does fine-tuning on all 10 datasets improve cross-dataset generalization?
% - Compare to Feger et al.'s joint benchmark training results for encoders
% - Leave-one-out evaluation: train on 9, test on 1

\subsection{Alternative Training Strategies}
\label{sec:ft_alternatives}

% - Optional extension, if fine-tuning degrades delta:
%   - Data augmentation with manipulated examples (train on M0 + M1 + M2)
%   - Early stopping based on manipulation robustness (not just F1)
%   - Curriculum learning: start with manipulated, then original
% - Goal: can we preserve robustness while gaining fine-tuning performance?
